% ===========================================================================
% CAPÍTULO 3: DESARROLLO
% ===========================================================================

\section{Preparación del entorno de desarrollo}

Para llevar a cabo la implementación propuesta en este Trabajo de Fin de Máster, fue
necesario establecer un entorno de desarrollo robusto que permitiera el codiseño
hardware-software (PS-PL) sobre la plataforma Zynq UltraScale+ MPSoC. A continuación, se
detalla la metodología seguida para la configuración de las herramientas de síntesis, la
compilación del sistema operativo en tiempo real (RTEMS) y la validación del flujo de trabajo
completo. El entorno de desarrollo principal se desplegó sobre una máquina virtual con una
distribución Linux (Ubuntu/Debian). En este sistema se instalaron las dependencias necesarias
para la compilación cruzada, incluyendo paquetes esenciales de desarrollo de C/C++,
Python 3 y herramientas de control de versiones.

\subsection{Instalación de Vivado y Vitis}

Para el desarrollo de la lógica programable (PL) y la generación de imágenes de arranque, se
utilizaron las herramientas oficiales de AMD Xilinx: Vivado Design Suite y Vitis Unified
Software Platform. Vivado se empleó para el diseño a nivel de hardware y la generación del
\textit{bitstream}, mientras que Vitis se utilizó para empaquetar los binarios y generar la imagen de
arranque (BOOT.bin).

\subsection{Instalación de RTEMS}

En lugar de utilizar binarios precompilados, se optó por compilar el sistema operativo RTEMS
(versión 7) desde cero, garantizando el control total sobre la configuración del núcleo. La
compilación se realizó utilizando la herramienta RTEMS Source Builder (RSB). Se generó un
\textit{toolchain} cruzado específico para la arquitectura AArch64. Se configuró y compiló el Board
Support Package (BSP) correspondiente a la plataforma ZynqMP (\texttt{zynqmp\_apu}),
habilitando explícitamente el soporte para multiprocesamiento simétrico (SMP).

Para validar el entorno recién compilado, se ejecutaron pruebas preliminares utilizando el
emulador QEMU (\texttt{qemu-system-aarch64}) configurado con la máquina
\texttt{xlnx-zcu102}. Se verificó la correcta ejecución de ejemplos básicos como
\texttt{hello.exe} y pruebas de rendimiento computacional como el \textit{benchmark} Dhrystone.

\subsection{Flujo de desarrollo completo}

Una vez preparadas todas las herramientas de desarrollo, se probó a realizar un desarrollo
completo de principio a fin implementando un ejemplo sencillo: una puerta AND. Para ello, se
siguieron los siguientes pasos.

\textbf{Parte Vivado.} Se creó un proyecto nuevo en Vivado llamado \texttt{and\_gate\_ps\_p}, y se
diseñó la puerta AND en VHDL:

\begin{lstlisting}[language=VHDL, caption={AND\_GATE.vhd --- ejemplo de puerta AND en VHDL}]
library IEEE;
use IEEE.STD_LOGIC_1164.ALL;

entity AND_GATE is
    Port ( A_B_IN : in  STD_LOGIC_VECTOR (1 downto 0);
           Z_OUT  : out STD_LOGIC);
end AND_GATE;

architecture Behavioral of AND_GATE is
begin
    Z_OUT <= A_B_IN(1) and A_B_IN(0);
end Behavioral;
\end{lstlisting}

Después se creó un \textit{Block Design}, se añadió el bloque Zynq UltraScale+ MPSoC IP y se
configuró automáticamente usando \textit{Block Automation} de Vivado. Se añadió el fichero VHDL
del paso anterior al \textit{Block Design}. Se añadieron dos AXI GPIO IPs al \textit{block design} para
comunicar el PS con la PL, y se conectaron los bloques entre sí, utilizando \textit{Automate
Connection} de Vivado para todas las conexiones faltantes.

Después se creó un \textit{Wrapper} del \textit{Block Design} para que Vivado pudiera sintetizarlo
correctamente. En la pestaña \textit{Address Editor} pueden verse las direcciones de los registros
para acceder a la parte PL desde el PS. Por último, se generó el \textit{BitStream} y se exportó el
hardware como .xsa.

\textbf{Generación de la imagen de arranque (Vitis).} Se siguieron los siguientes pasos para
empaquetar el hardware generado en Vivado en un binario que el MPSoC pudiera ejecutar:
se creó un nuevo componente en Vitis de tipo \textit{Platform}, utilizando como base el .xsa
generado; se compiló el FSBL usando el ejemplo \textit{Zynq MP FSBL}; y se generó el archivo
\texttt{BOOT.bin} mediante el generador de imágenes de Vitis con la estructura de componentes
estándar (FSBL, PMUFW, bitstream de la PL, BL31, U-Boot y DTB).

\textbf{Parte RTEMS.} Una vez generado el archivo de arranque, se desarrolló un software de
ejemplo que interactuaba con los puertos GPIO conectados por AXI, excitando todas las
combinaciones posibles de la puerta AND en la PL y leyendo su salida. Fue necesario añadir
un archivo de mapeado de memoria para que el sistema operativo reconociera las direcciones
correspondientes a la PL, y un \texttt{wscript} para la compilación.

Para facilitar el desarrollo rápido en la parte de RTEMS, se desarrolló un script que compila
la imagen de RTEMS, generando \texttt{rtems.img}. Para evitar quitar y poner la tarjeta SD
constantemente, se desarrolló un script en Python que es capaz de subir la imagen de RTEMS
a la tarjeta SD mediante una conexión USB.

Una vez verificado que esta metodología de desarrollo funcionaba correctamente, se
estableció como el proceso estándar a seguir para los próximos desarrollos.

% ---------------------------------------------------------------------------
\section{Diseño del transceptor serie configurable (PL)}
% ---------------------------------------------------------------------------

El transceptor serie se diseñó para ser altamente configurable y adaptable a cualquier
dispositivo. Los parámetros configurables son:

\begin{itemize}
    \item Baudrate (desde 50 hasta 4.000.000 baudios).
    \item Número de \textit{stop bits} (1, 1,5 o 2).
    \item Paridad (par, impar, marca, espacio o deshabilitada).
    \item Orden de los bits (LSB o MSB).
    \item Número de bits de datos (entre 5 y 9).
\end{itemize}

\subsection{Transmisor}

El siguiente diagrama FSM resume el comportamiento del transmisor. Debido a su tamaño, se
ha dividido en tres partes para su legibilidad.

\begin{figure}[H]
    \centering
    \fbox{\parbox{0.85\textwidth}{\centering \textit{[Figura pendiente: Diagrama FSM del transmisor (tres partes)]}\vspace{0.5cm}}}
    \caption{Diagrama FSM del transmisor serie.}
    \label{fig:fsm_tx}
\end{figure}

Durante las pruebas sobre la placa CDHS se detectó que en las primeras recepciones aparecían
ocasionalmente caracteres corruptos. Como primera hipótesis se consideró que la señal DE
podría estar conmutando demasiado rápido al paso de transmisión a recepción, haciendo que
el transceptor físico no tuviera tiempo de establecerse. Para corregirlo, se añadieron dos
estados adicionales en la FSM del transmisor que introducen un retardo de microsegundos
entre la desactivación del último bit y la desactivación de DE. Sin embargo, esto no resolvió
el problema, descartando la hipótesis inicial.

\subsection{Receptor}

El siguiente diagrama FSM resume el comportamiento del receptor.

\begin{figure}[H]
    \centering
    \fbox{\parbox{0.85\textwidth}{\centering \textit{[Figura pendiente: Diagrama FSM del receptor serie]}\vspace{0.5cm}}}
    \caption{Diagrama FSM del receptor serie.}
    \label{fig:fsm_rx}
\end{figure}

\subsubsection{Decisión de diseño: verificación de start-bit a mitad de período}

Tras descartar el flanco de DE como causa de los caracteres corruptos, el análisis se centró
en el receptor. El protocolo serie comienza con un \textbf{start-bit}: la línea, normalmente en
reposo en nivel alto, cae a nivel bajo para señalizar el inicio de una trama. El receptor detecta
este flanco descendente y arranca la máquina de estados de recepción.

El problema era que pulsos de ruido breves en la línea generaban flancos descendentes
espurios que el receptor interpretaba como \textit{start-bits} válidos, capturando una trama de
basura y enviando un carácter corrupto al \textit{buffer}.

La solución se implementó directamente en la FSM: en el estado de recepción del
\textit{start-bit}, una vez transcurrido el \textbf{primer medio período de bit}, se vuelve a muestrear la
línea antes de continuar. Si la línea ha vuelto a nivel alto, el pulso era ruido y la FSM regresa
al estado IDLE sin capturar nada. Solo si la línea sigue en bajo al llegar a la mitad del bit se
considera un \textit{start-bit} legítimo y se procede con la recepción del resto de la trama.

Esta técnica ---verificación a mitad de \textit{start-bit}--- es una práctica estándar de robustez en
receptores UART. Su implementación en la FSM VHDL eliminó por completo la aparición de
caracteres corruptos en las pruebas posteriores.

\subsection{Oscilador controlado numéricamente (NCO)}

Este oscilador sirve para generar una señal de \textit{enable} a un baudrate seleccionado. Hay 54
valores de velocidad seleccionables almacenados en una tabla de verdad. Para poder generar
frecuencias que no son divisibles exactamente de la frecuencia de reloj de la FPGA, se utiliza
este circuito que corrige periódicamente el desfase generado por esta imprecisión, logrando
replicar un amplio número de frecuencias con muy poco error.

El circuito consiste en un sumador acumulador: para cada frecuencia seleccionada, se tiene
calculado un incremento concreto que depende del número de bits del sumador. Cuantos más
bits tenga este sumador, mejor precisión se consigue. En esta implementación se han
utilizado 32 bits. El incremento se suma y acumula en cada período de reloj; cuando el
sumador se desborda, se utiliza el bit de \textit{carry\_out} como \textit{tick} del NCO. Este \textit{tick} se genera
a la frecuencia deseada, y el mecanismo de acumulación garantiza que el error promedio sea
inferior al de un divisor entero simple. La tabla de frecuencias soportadas y su error medido
se recoge en el Anexo~1.

\begin{figure}[H]
    \centering
    \fbox{\parbox{0.85\textwidth}{\centering \textit{[Figura pendiente: Diagrama del circuito NCO de 32 bits]}\vspace{0.5cm}}}
    \caption{Circuito del oscilador controlado numéricamente (NCO).}
    \label{fig:nco}
\end{figure}

\subsection{Shift-register}

El registro de desplazamiento está diseñado para poder almacenar datos tanto si la
configuración es LSB-\textit{first} como MSB-\textit{first}, de modo que una vez recibido el mensaje la
palabra en la salida siempre es la correcta independientemente del modo configurado.

\subsection{FIFO}

La FIFO es un bloque IP con 9 bits de datos (para poder almacenar mensajes de 9 bits) y
512 de profundidad, que permite almacenar temporalmente varios mensajes y habilitarlos
para su lectura posterior por el PS.

\subsection{Bloque Zynq UltraScale+ MPSoC}

Para que el proyecto en Vivado funcione correctamente, es importante aplicar el \textit{preset} al
bloque IP Zynq UltraScale+ antes de añadir el resto de elementos, ya que se aplican
configuraciones de relojes y alimentación cruciales para el correcto funcionamiento de la PL.

% ---------------------------------------------------------------------------
\section{Herramientas para facilitar el desarrollo del transceptor (PL)}
% ---------------------------------------------------------------------------

\subsection{Generación de transceivers con TCL}

Se desarrolló un script de TCL que permite generar y conectar automáticamente un número
deseado de transceptores en el \textit{Block Design} de Vivado. Para que funcione correctamente, es
necesario tener las fuentes .vhd del transceptor en el proyecto, el \textit{Block Design} creado, y el
bloque IP Zynq UltraScale+ MPSoC importado y configurado previamente; a partir de ahí, el
script se encarga de añadir y conectar el resto de los bloques necesarios.

\subsection{Generación de proyecto con TCL}

También se desarrolló otro script de TCL (\texttt{regenerate\_all.tcl}) que regenera el proyecto
entero de Vivado con un número determinado de transceptores, partiendo desde cero. Esto
permite reproducir exactamente cualquier configuración sin necesidad de guardar el proyecto
de Vivado completo en el repositorio.

% ---------------------------------------------------------------------------
\section{Generación del binario en Vitis (PS)}
% ---------------------------------------------------------------------------

Una vez diseñada la lógica hardware en Vivado y exportado el hardware, se siguieron los
pasos establecidos previamente para generar el binario BOOT.bin.

El proceso manual de exportación del XSA, creación de la plataforma Vitis, compilación del
FSBL y ensamblado del BOOT.BIN con \texttt{bootgen} resultó tedioso de repetir en cada
iteración del diseño. Para automatizarlo, se desarrolló el script \texttt{generate\_boot.sh}
(disponible en \texttt{tfm/04\_tools/}), un \textit{wizard} interactivo en Bash que encadena todos
estos pasos de forma desatendida. El script ofrece tres puntos de entrada según el estado en
que se encuentre el trabajo: partir del proyecto de Vivado y generar \textit{bitstream} desde cero,
partir de un proyecto con \textit{bitstream} ya generado y solo exportar el XSA, o partir
directamente de un XSA ya exportado. A partir de ahí invoca Vitis en modo script mediante
su API Python para crear la plataforma y compilar el FSBL, y llama a \texttt{bootgen} para
ensamblar el BOOT.BIN final con los componentes de arranque precompilados (PMUFW,
BL31, U-Boot, DTB). Esto redujo el tiempo de iteración hardware $\rightarrow$ imagen de varios
minutos de clics a una sola ejecución de terminal.

% ---------------------------------------------------------------------------
\section{Diseño del driver serie en RTEMS (PS)}
% ---------------------------------------------------------------------------

\subsection{Motivación y contexto}

El subsistema de comunicaciones del OBC requiere gestionar simultáneamente múltiples
interfaces serie físicas (RS422 y RS485) desde una aplicación de tiempo real ejecutada sobre
RTEMS. El bloque hardware \texttt{CONFIGURABLE\_SERIAL\_TOP}, implementado en VHDL en
la lógica programable (PL), expone cada transceptor a través de un periférico AXI GPIO de
doble canal, cuya interfaz de bajo nivel es demasiado específica para ser usada directamente
desde la aplicación. Para abstraer esta complejidad y ofrecer una API uniforme e
independiente de la dirección física de cada instancia, se desarrollaron los ficheros
\texttt{transceiver.c} y \texttt{transceiver.h}.

El diseño cubre tres requisitos fundamentales: (1) descubrimiento dinámico del hardware en
tiempo de arranque, de modo que el mismo binario funcione con cualquier número de
transceptores instanciados en el \textit{bitstream}; (2) comunicación orientada a interrupciones
para no bloquear el procesador durante la espera de datos; y (3) una API sencilla que
permita a la aplicación enviar y recibir datos sin conocer los detalles del mapa de registros.

\subsection{Interfaz hardware: mapa de registros}

Cada instancia del transceptor ocupa una región de 4~KB del espacio de memoria del
procesador, estructurada como un AXI GPIO de doble canal:

\begin{itemize}
    \item \textbf{Canal 1 (escritura, PS $\rightarrow$ PL):} contiene el dato de transmisión (bits 8:0),
    los bits de control de protocolo (SEND, ERROR\_OK, DATA\_READ) y los campos de
    configuración: tasa de baudios (6 bits), bits de parada (2 bits), paridad (3 bits), bits de
    datos (3 bits), orden de bits y modo SLO.

    \item \textbf{Canal 2 (lectura, PL $\rightarrow$ PS):} contiene el byte recibido (bits 8:0) y las
    banderas de estado: RX\_EMPTY, RX\_FULL, PARITY\_ERROR, FRAME\_ERROR y TX\_RDY.
\end{itemize}

Adicionalmente, un bloque de información del sistema situado en la dirección fija
\texttt{0xA0020000} expone en tiempo de ejecución el número de transceptores instanciados, el
\textit{stride} de memoria entre instancias y la dirección base de la primera de ellas. Más allá del
último transceptor se aloja el controlador de interrupciones AXI INTC, compartido por todas
las instancias.

\subsection{Arquitectura software}

La librería se organiza en torno a dos estructuras de datos centrales:

\textbf{\texttt{Transceiver\_Config\_t}} agrupa los parámetros de configuración del protocolo: tasa de
baudios, número de bits de datos (5 a 9), paridad (par, impar, \textit{mark}, \textit{space} o ninguna),
bits de parada (1, 1,5 o 2), orden de bits (LSB o MSB \textit{first}) y modo de emisión lenta (SLO,
que reduce las emisiones electromagnéticas). Todos los valores se expresan mediante macros
simbólicas definidas en \texttt{transceiver.h}, que corresponden directamente a los valores
codificados en la LUT del NCO hardware.

\textbf{\texttt{Transceiver}} es el objeto \textit{handle} que representa una instancia concreta del
transceptor en ejecución. Contiene el mapa de direcciones calculado dinámicamente
(\texttt{base\_addr}, \texttt{intc\_base} y los \textit{offsets} derivados), las máscaras de interrupción
individuales de RX y TX, y el estado software: un \textit{buffer} circular de recepción y uno de
transmisión (ambos de 4096 bytes por defecto, asignados dinámicamente), los
identificadores de las primitivas RTEMS asociadas (tarea \textit{worker}, \textit{mutex} de RX y \textit{mutex} de
TX) y el \textit{callback} de usuario para notificación de datos recibidos.

\subsection{Descubrimiento dinámico del hardware}

La función \texttt{Transceiver\_Hardware\_Discover()} lee el bloque de información del sistema al
arranque y extrae el número de instancias presentes, el \textit{stride} de memoria y la dirección
base. Con estos tres valores se calcula automáticamente la dirección del INTC compartido:

\begin{center}
\texttt{g\_intc\_addr = g\_hw\_base + (g\_hw\_count $\times$ g\_hw\_stride)}
\end{center}

Este mecanismo permite que el mismo binario RTEMS funcione sin modificación con cualquier
número de transceptores (hasta el máximo soportado de 14), cubriendo configuraciones que
van desde un único canal hasta la instalación completa. La alternativa de \textit{hardcodear} las
direcciones habría requerido recompilar la aplicación para cada variante del \textit{bitstream}.

\subsection{Modelo de interrupciones: ISR maestra y tareas worker}

El diseño de la gestión de interrupciones adopta el patrón \textit{deferred interrupt processing}
recomendado para RTEMS. Se instala una única ISR en el vector 121 (\texttt{Master\_ISR}) que
atiende a todos los transceptores a través del AXI INTC compartido. La ISR funciona del
siguiente modo:

\begin{enumerate}
    \item Lee el registro de interrupciones pendientes (INTC\_ISR).
    \item Para cada transceptor con interrupción RX activa: reconoce la interrupción
    (INTC\_IAR), deshabilita temporalmente la línea (INTC\_CIE) y envía el evento
    \texttt{RTEMS\_EVENT\_0} a la tarea \textit{worker} correspondiente.
    \item Para cada transceptor con interrupción TX activa: extrae el siguiente byte del
    \textit{buffer} circular de transmisión, lo escribe en el registro de control pulsando
    \texttt{BIT\_TX\_SEND}, y limpia la interrupción. Si el \textit{buffer} ha quedado vacío, marca
    \texttt{tx\_busy = false}.
\end{enumerate}

Cada instancia \texttt{Transceiver} tiene asociada una tarea RTEMS dedicada
(\texttt{Rx\_Worker\_Task}) que se bloquea esperando el evento. Al recibirlo, vacía el FIFO
hardware byte a byte hacia el \textit{buffer} circular de recepción, pulsando \texttt{BIT\_DATA\_READ}
para confirmar la lectura de cada byte al hardware. El acceso al \textit{buffer} se protege con un
semáforo binario de prioridad para evitar condiciones de carrera. Tras drenar el FIFO, la
\textit{worker} re-habilita la interrupción (INTC\_SIE) e invoca el \textit{callback} de usuario si está
registrado.

La separación entre la ISR (que solo despierta la tarea) y la \textit{worker} (que realiza el trabajo
real) garantiza que el tiempo de latencia en contexto de interrupción sea mínimo y que el
procesamiento se ejecute con las prioridades del planificador RTEMS.

\subsection{Inicialización y configuración del hardware}

\texttt{Transceiver\_Init()} realiza la puesta en marcha completa de una instancia: calcula la
dirección base y la dirección del INTC, asigna las máscaras de interrupción individuales (bit
$2 \cdot id$ para RX y bit $2 \cdot id + 1$ para TX dentro del registro de 32 bits del INTC),
asigna dinámicamente los \textit{buffers} circulares, crea el semáforo y la tarea \textit{worker} con
\texttt{rtems\_semaphore\_create} y \texttt{rtems\_task\_create}, escribe la configuración de
protocolo en el hardware y lanza la tarea \textit{worker}.

La función \texttt{Transceiver\_Global\_INIT()} orquesta la inicialización completa del sistema:
primero llama a \texttt{Transceiver\_Hardware\_Discover()} para poblar las variables globales, y a
continuación a \texttt{Transceiver\_Global\_INTC\_Init()}, que habilita el modo maestro del INTC
(registro MER) e instala la \texttt{Master\_ISR}.

\subsection{API pública}

La interfaz expuesta al código de aplicación se reduce a cinco funciones:

\begin{table}[H]
    \centering
    \begin{tabular}{ll}
        \hline
        \textbf{Función} & \textbf{Descripción} \\
        \hline
        \texttt{Transceiver\_Global\_INIT()} & Descubrimiento hardware e inicialización del INTC. \\
        \texttt{Transceiver\_Init(dev, id, cfg)} & Inicializa la instancia \texttt{dev} con el id y configuración. \\
        \texttt{Transceiver\_SetRxCallback(dev, cb, arg)} & Registra \textit{callback} de recepción. \\
        \texttt{Transceiver\_Read(dev, buf, maxlen)} & Extrae hasta \texttt{maxlen} bytes del buffer RX. \\
        \texttt{Transceiver\_SendString(dev, s)} & Encola la cadena \texttt{s} y arranca la transmisión. \\
        \hline
    \end{tabular}
    \caption{API pública del driver serie.}
    \label{tab:api}
\end{table}

\subsection{Decisiones de diseño destacadas}

\textbf{Buffer circular frente a copia directa en ISR.} Dado que RTEMS no permite operaciones de
copia de longitud arbitraria en contexto de interrupción sin afectar la latencia del sistema, se
optó por \textit{buffers} circulares gestionados desde la tarea \textit{worker}. El tamaño de 4096 bytes
proporciona margen suficiente para ráfagas de datos a las tasas más altas soportadas (hasta
4~Mbaudios).

\textbf{Transmisión dirigida por interrupción.} El primer byte de cada cadena a transmitir se escribe
directamente desde \texttt{Transceiver\_SendString()} para arrancar el hardware; los bytes
siguientes son enviados sucesivamente por la \texttt{Master\_ISR} al recibir la interrupción
TX\_RDY, sin necesidad de que la tarea de aplicación permanezca bloqueada durante la
transmisión.

\textbf{Compatibilidad multi-instancia mediante tabla global.} El array \texttt{g\_instances[]} permite
a la \texttt{Master\_ISR} localizar en $O(1)$ el objeto \texttt{Transceiver} asociado a cada par de bits de
interrupción, evitando búsquedas costosas en contexto de ISR.

\textbf{Modo SLO.} El bit \texttt{BIT\_SLO} del registro de control activa en el hardware VHDL una
rampa de subida más lenta en las señales de salida, reduciendo el contenido espectral de alta
frecuencia y mejorando la compatibilidad electromagnética en entornos con restricciones EMI.

% ---------------------------------------------------------------------------
\section{Diseño hardware}
% ---------------------------------------------------------------------------

Una vez completado el firmware del transceptor serie, se desarrolló el hardware de prueba
necesario para validar el sistema en condiciones reales. El plan inicial era diseñar una única
PCB propia para testear múltiples líneas serie. Posteriormente, el proyecto LINCE requirió
dos tarjetas adicionales con especificaciones impuestas por Indra ---la placa CDHS y la placa
AOCS--- para que Sener pudiera validar su firmware sobre la ZCU102. Los esquemáticos
completos, BOMs y archivos de fabricación de las tres placas están disponibles en el
repositorio del proyecto bajo \texttt{HARDWARE/}.

\subsection{Diseño de la placa de comunicación serie (diseño propio)}

El objetivo de esta placa es poder ejercitar simultáneamente los 14 transceptores del IP core,
replicando las topologías de bus reales que se encontrarán en el satélite: RS485 multipunto
con múltiples nodos en el mismo cable, y RS422 en configuración \textit{master/slave} con cruce de
TX y RX. A diferencia de las placas CDHS y AOCS, cuyas especificaciones vinieron dictadas
por Indra, esta placa fue diseñada con total libertad para maximizar la flexibilidad de los
ensayos en laboratorio.

La arquitectura general de la placa divide los 14 canales en dos grupos conectados al FMC:

\begin{itemize}
    \item \textbf{7 drivers RS485 (RS0--RS6)} conectados a un bus diferencial común mediante
    \textit{jumpers}, formando una topología multipunto configurable.
    \item \textbf{7 drivers RS422 (RS7--RS13)} organizados en dos buses \textit{master/slave}
    independientes (BUS 1: 1 \textit{master} + 3 \textit{slaves}; BUS 2: 1 \textit{master} + 2 \textit{slaves}), con la línea
    TX cruzada con RX para emular la conexión real punto a punto RS422.
\end{itemize}

\begin{figure}[H]
    \centering
    \fbox{\parbox{0.85\textwidth}{\centering \textit{[Figura pendiente: arquitectura general de la placa --- 7 drivers RS485, conector FMC, 7 drivers RS422]}\vspace{0.5cm}}}
    \caption{Arquitectura general de la placa de comunicación serie.}
    \label{fig:placa_serial_top}
\end{figure}

\subsubsection{Topología RS485: bus compartido configurable por jumper}

La decisión de diseño más relevante de esta placa es el mecanismo de bus compartido RS485
sin necesidad de recablear. Cada driver RS485 tiene un conector Dupont de 3 pines
(A+, B$-$, GND) que actúa a la vez como punto de conexión hacia el exterior y como punto de
interconexión entre drivers adyacentes. Colocando un \textit{jumper} entre dos conectores
consecutivos, las líneas A y B de ambos drivers quedan físicamente unidas, formando un
segmento de bus RS485 compartido. Retirando el \textit{jumper}, el driver queda independiente y
puede conectarse a un periférico externo de forma individual.

Este mecanismo permite configurar por hardware cualquier partición de los 7 drivers RS485
en uno o varios buses, sin modificar el firmware ni el cableado exterior.

\subsubsection{Selector de conector Micro-D para RS485}

Los conectores de mayor densidad (Micro-D) de la placa son compartidos entre el Driver 0 y
el Driver 6 mediante un \textit{jumper} de selección de hardware. Dependiendo de la posición del
\textit{jumper}, los conectores Micro-D quedan asignados a uno u otro driver, lo que permite usar
un único tipo de conector con cableado de satélite sin duplicar el número de conectores en
la placa.

\subsubsection{Topología RS422: master/slave configurable por jumper}

Para RS422, cada driver \textit{slave} dispone de un conector con las líneas TX cruzadas a RX
(TX-Y+/TX-Z$-$ conectadas a RX del bus), replicando la conexión estándar punto a punto
RS422 donde el receptor del esclavo escucha las transmisiones del \textit{master}. Al igual que en
RS485, un \textit{jumper} determina si el driver se une al bus común o permanece independiente.

El esquemático completo (12 hojas) y la BOM están disponibles en
\texttt{HARDWARE/lince\_comunicacion\_serial/} del repositorio del proyecto.

% ----- CDHS -----
\subsection{Diseño de la placa CDHS}

La placa LINCE3 CDHS BreakoutBox es una tarjeta de interconexión entre la ZCU102
(conectada por FMC) y los periféricos del subsistema de gestión de datos y calentadores del
satélite. Sus especificaciones fueron definidas por Indra: 6 conectores D-Sub-9 (J1--J6),
interfaces CAN redundante, tres canales RS422/RS485, cuatro líneas PWM para calentadores
y un ADC para termistores.

\begin{figure}[H]
    \centering
    \fbox{\parbox{0.85\textwidth}{\centering \textit{[Figura pendiente: diagrama jerárquico CDHS --- bloque CAN (J1), RS (J2--J4), PWM (J5), ADC/THERM (J6)]}\vspace{0.5cm}}}
    \caption{Diagrama de bloques de la placa CDHS.}
    \label{fig:cdhs_top}
\end{figure}

\begin{figure}[H]
    \centering
    \fbox{\parbox{0.85\textwidth}{\centering \textit{[Figura pendiente: render 3D de la PCB CDHS]}\vspace{0.5cm}}}
    \caption{Render 3D de la placa CDHS.}
    \label{fig:cdhs_3d}
\end{figure}

\subsubsection{Decisión de nivel de tensión: 1,8 V}

Todos los bancos de pines del conector FMC HPC utilizados por esta placa operan a 1,8~V.
Se buscaron deliberadamente componentes con VIO nativo a 1,8~V para evitar etapas de
adaptación de nivel adicionales entre el MPSoC y los transceptores. Esta decisión simplifica
el diseño y reduce el número de componentes activos en la cadena de señal.

\subsubsection{Subsistema CAN}

El bus CAN implementa topología redundante con dos instancias independientes: CAN
Nominal (CAN\_NOM) y CAN Redundante (CAN\_RED). Se seleccionó el transceptor
\textbf{TCAN1044AVDRQ1} de Texas Instruments por su compatibilidad nativa con VIO de 1,8~V,
grado automotriz y baja corriente en modo \textit{standby}. Ambos canales convergen en el
conector J1.

Las resistencias de terminación siguen el esquema \textit{split termination} recomendado en las
notas de aplicación del estándar CAN (ISO 11898-2). En lugar de una única resistencia de
120~$\Omega$ entre CANH y CANL, se emplean \textbf{dos resistencias de 60~$\Omega$ en serie} con un
\textbf{condensador de 4,7~nF} conectado desde el punto medio al plano de masa. Cada una de
las dos resistencias de 60~$\Omega$ está controlada por un \textit{jumper} independiente, lo que
permite habilitar o deshabilitar la terminación sin modificar el circuito. La resistencia
equivalente sigue siendo 120~$\Omega$ cuando ambos \textit{jumpers} están colocados, y el
condensador de derivación crea un camino de baja impedancia para las interferencias de
modo común de alta frecuencia hacia masa, mejorando el comportamiento EMC del bus.

Las líneas CANH y CANL de cada canal incorporan diodos de protección ESD
\textbf{ESDCAN24-2BLY}, específicamente diseñados para buses CAN, que clampean transitorios
por encima de su tensión de ruptura sin introducir una capacitancia parásita significativa.

\begin{figure}[H]
    \centering
    \fbox{\parbox{0.85\textwidth}{\centering \textit{[Figura pendiente: esquemático CAN.SchDoc --- transceptores TCAN1044 con terminación \textit{split} y diodos ESD]}\vspace{0.5cm}}}
    \caption{Subsistema CAN de la placa CDHS.}
    \label{fig:cdhs_can}
\end{figure}

\subsubsection{Subsistema RS422/RS485}

Se disponen tres canales serie idénticos e independientes (RS1, RS2, RS3), cada uno con el
transceptor \textbf{THVD1424RGTR}. Este integrado fue elegido porque soporta VIO de 1,8~V,
incorpora resistencias de terminación internas conmutables y permite seleccionar entre modo
\textit{Half-Duplex} y \textit{Full-Duplex} sin cambiar el hardware. La configuración se expone mediante
\textit{jumpers} físicos de 2 pines: \textbf{H/F} (conmuta entre \textit{Half-Duplex} y \textit{Full-Duplex}), \textbf{SLR}
(habilita control de \textit{slew rate}), \textbf{TERM\_TX / TERM\_RX} (conectan las resistencias de
terminación internas). Los tres canales salen por los conectores J2, J3 y J4.

Tanto las líneas RX como las TX de cada canal llevan un par de diodos TVS
\textbf{SM712-02HTG} conectados entre las líneas diferenciales y el plano de masa.

\paragraph{Decisión de diseño: cortocircuito DE--RE.}
El THVD1424 expone dos pines de control independientes: \textbf{DE} (\textit{Driver Enable}, activo en
alto) habilita el transmisor, y \textbf{RE} (\textit{Receiver Enable}, activo en bajo) habilita el receptor.
En el esquemático, ambos pines están conectados a la misma señal de control procedente del
MPSoC, de modo que una única línea digital controla simultáneamente transmisor y receptor
de forma complementaria (tabla~\ref{tab:de_re}).

\begin{table}[H]
    \centering
    \begin{tabular}{ccc}
        \hline
        \textbf{Señal DE/RE} & \textbf{Transmisor (DE)} & \textbf{Receptor (RE activo bajo)} \\
        \hline
        1 (alto) & Habilitado & Deshabilitado \\
        0 (bajo) & Deshabilitado & Habilitado \\
        \hline
    \end{tabular}
    \caption{Control complementario DE/RE en el THVD1424.}
    \label{tab:de_re}
\end{table}

Esta decisión responde a dos motivos. Por un lado, el equipo de Indra comunicó que en la
arquitectura del satélite los esclavos solo transmiten bajo demanda del OBC, por lo que no
existe un escenario real en el que un nodo deba transmitir y recibir simultáneamente. Por
otro lado, en modo \textit{Half-Duplex} las líneas TX y RX comparten físicamente el mismo par
diferencial A/B; si el receptor permaneciera activo durante la transmisión, el nodo escucharía
su propio eco.

\begin{figure}[H]
    \centering
    \fbox{\parbox{0.85\textwidth}{\centering \textit{[Figura pendiente: RS.SchDoc --- circuito completo de un canal RS con THVD1424RGTR y diodos TVS]}\vspace{0.5cm}}}
    \caption{Canal RS de la placa CDHS con THVD1424RGTR.}
    \label{fig:cdhs_rs}
\end{figure}

\subsubsection{Subsistema PWM (calentadores)}

Cuatro señales PWM de 1,8~V procedentes del FMC se elevan a 3,3~V mediante el adaptador
de niveles \textbf{TXU0104PWR} para excitar los circuitos de control de calentadores externos.
Las cuatro líneas adaptadas salen por J5.

Para la validación de esta interfaz se diseñó un bloque VHDL autónomo,
\texttt{PWMx4\_auto\_test}, sintetizado en la PL junto al resto del diseño de Vivado. El bloque
genera cuatro señales PWM independientes con \textit{duty cycle} del 50\% a partir del reloj de
100~MHz de la FPGA: canal 0 a 10~kHz, canal 1 a 5~kHz, canal 2 a 1~kHz y canal 3 a
100~Hz. Al ser completamente autónomo no requiere ningún driver ni intervención del PS, lo
que permitió verificar el subsistema PWM de la placa con el mismo BOOT.BIN utilizado para
las pruebas serie.

\subsubsection{Subsistema ADC (termistores)}

Cuatro entradas analógicas de termistores (CH0--CH3) se digitalizan con el ADC SAR de
12 bits \textbf{ADS7950QDBTRQ1}, comunicado al procesador por SPI a 1,8~V. La circuitería
analógica opera a 3,3~V; ambos carriles (+VBD a 1,8~V y +VA a 3,3~V) son independientes
para evitar acoplos entre el dominio digital y el analógico \cite{ti2018ads7950}. Los
termistores se conectan por J6.

\subsubsection{Alimentación}

El conector FMC suministra 12~V, 3,3~V y el carril VADJ a 1,8~V. El convertidor conmutado
\textbf{R-78E5.0-1.0} genera los 5~V necesarios para la etapa de potencia de los transceptores
RS y CAN a partir de los 12~V del FMC. Se eligió este módulo DC/DC por su integración en
un \textit{footprint} reducido de 3 pines (equivalente a un regulador lineal) y su eficiencia $\geq
96\%$, evitando la disipación térmica que tendría un LDO con ese diferencial de tensión.

\subsubsection{Conectores externos}

Los seis puertos D-Sub-9 (J1--J6) son de tipo macho o hembra según la convención de uso.
Los escudos metálicos de todos los conectores están conectados al plano de GND para
mantener la continuidad de apantallamiento EMI con los cables blindados.

El mapa completo de señales entre la placa CDHS y el conector FMC se encuentra en el
Anexo~1 de este documento. El esquemático completo y la BOM están disponibles en
\texttt{HARDWARE/CDHS/}.

% ----- AOCS -----
\subsection{Diseño de la placa AOCS}

La placa LINCE3 AOCS BreakoutBox es la tarjeta de interconexión para el subsistema de
control de actitud y órbita del satélite. Al igual que la CDHS, conecta la ZCU102 mediante
FMC y expone las interfaces al exterior a través de 8 conectores D-Sub-9 (J1--J8).

\begin{figure}[H]
    \centering
    \fbox{\parbox{0.85\textwidth}{\centering \textit{[Figura pendiente: diagrama jerárquico AOCS --- bloques RS (J1, J3--J5, J8), MOT-PWM (J2), SpaceWire (J6, J7)]}\vspace{0.5cm}}}
    \caption{Diagrama de bloques de la placa AOCS.}
    \label{fig:aocs_top}
\end{figure}

\begin{figure}[H]
    \centering
    \fbox{\parbox{0.85\textwidth}{\centering \textit{[Figura pendiente: render 3D de la PCB AOCS]}\vspace{0.5cm}}}
    \caption{Render 3D de la placa AOCS.}
    \label{fig:aocs_3d}
\end{figure}

\subsubsection{Canales RS422/RS485}

La placa incorpora 5 canales serie (RS1, RS3, RS4, RS5, RS8) con exactamente el mismo
diseño de transceptor THVD1424RGTR descrito en la sección de la placa CDHS, incluyendo
el cortocircuito DE--RE y los \textit{jumpers} de configuración H/F, SLR y terminación.

\subsubsection{Canales SpaceWire}

Se implementan dos enlaces SpaceWire \textit{full-duplex} independientes (SPW1 y SPW2) mediante
señalización diferencial LVDS, sin transceptor activo adicional, ya que el estándar SpaceWire
opera directamente a niveles LVDS compatibles con los bancos del FMC. Los pares
diferenciales ---cuatro por enlace: DIN, SIN, SOUT y DOUT--- se han trazado en la PCB con
técnicas de igualación de longitud (\textit{length matching}) para minimizar el \textit{skew} entre el par
de datos y el par de \textit{strobe}, y con impedancia característica controlada de 100~$\Omega$
diferencial.

\begin{figure}[H]
    \centering
    \fbox{\parbox{0.85\textwidth}{\centering \textit{[Figura pendiente: esquemático SpaceWire --- cuatro pares diferenciales y su conexión al FMC]}\vspace{0.5cm}}}
    \caption{Subsistema SpaceWire de la placa AOCS.}
    \label{fig:aocs_spw}
\end{figure}

\subsubsection{Control de motores (MOT-PWM)}

Seis señales PWM de 1,8~V procedentes del FMC se elevan al nivel de tensión requerido por
los puentes en H externos mediante un adaptador de niveles. Las señales se agrupan en tres
pares (X, Y, Z para los tres ejes del satélite), con dos líneas por eje para controlar la
dirección de giro del motor. Salen por J2. La alimentación de potencia de los motores se
conecta mediante dos bornes banana independientes de la alimentación de la placa.

\subsubsection{Alimentación}

Misma arquitectura que la placa CDHS: 12~V del FMC, convertidor DC/DC para los 5~V de
potencia, y VADJ 1,8~V para los transceptores.

El mapa completo de señales entre la placa AOCS y el conector FMC se encuentra en el
Anexo~1. El esquemático completo y la BOM están disponibles en \texttt{HARDWARE/AOCS/}.

% ----- FABRICACIÓN -----
\subsection{Fabricación}

Una vez validados los esquemáticos y completado el rutado de las PCBs, se encargaron las
placas con stencil a \textbf{PCBWay}, y los componentes a \textbf{Mouser} y \textbf{DigiKey} según
disponibilidad. El proceso de montaje en el laboratorio siguió los pasos habituales de
soldadura por reflujo SMD:

\begin{enumerate}
    \item \textbf{Aplicación de pasta de soldadura.} Se fijó la placa a la mesa con placas de
    idéntico grosor alrededor y cinta de pintor. Se colocó el stencil alineado y se extendió
    la pasta con espátula.
    \item \textbf{Colocación de componentes.} Con pinzas de punta fina, se colocaron los
    componentes SMD sobre la pasta. Se montaron dos placas en paralelo para optimizar el
    tiempo de proceso.
    \item \textbf{Reflujo.} Las placas se introdujeron en el horno de reflujo del laboratorio
    seleccionando el perfil de temperatura adecuado para la pasta de soldadura utilizada.
\end{enumerate}

\begin{figure}[H]
    \centering
    \fbox{\parbox{0.85\textwidth}{\centering \textit{[Figura pendiente: foto del proceso de fabricación --- aplicación de pasta con stencil]}\vspace{0.5cm}}}
    \caption{Aplicación de pasta de soldadura con stencil.}
    \label{fig:fab_pasta}
\end{figure}

\begin{figure}[H]
    \centering
    \fbox{\parbox{0.85\textwidth}{\centering \textit{[Figura pendiente: foto de la placa CDHS soldada (cara superior)]}\vspace{0.5cm}}}
    \caption{Placa CDHS tras el proceso de soldadura.}
    \label{fig:cdhs_soldada}
\end{figure}

\begin{figure}[H]
    \centering
    \fbox{\parbox{0.85\textwidth}{\centering \textit{[Figura pendiente: foto de la placa AOCS soldada (cara superior)]}\vspace{0.5cm}}}
    \caption{Placa AOCS tras el proceso de soldadura.}
    \label{fig:aocs_soldada}
\end{figure}

% ---------------------------------------------------------------------------
\section{Desarrollo de herramientas de testing}
% ---------------------------------------------------------------------------

\subsection{Aplicación de testing del driver serie en RTEMS}

La aplicación de testing del driver serie está estructurada en tres archivos:

\begin{itemize}
    \item \texttt{init.c}: configuración del sistema RTEMS mediante macros de
    \texttt{confdefs.h}.
    \item \texttt{transceiver.c} + \texttt{transceiver.h}: driver del transceptor serie sobre AXI GPIO.
    \item \texttt{main.c}: aplicación de testing que usa la API del driver.
\end{itemize}

\subsubsection{Configuración del sistema RTEMS --- \texttt{init.c}}

\texttt{init.c} es el fichero de configuración estática de RTEMS. No contiene lógica de
aplicación; define el entorno de ejecución mediante macros que \texttt{<rtems/confdefs.h>}
convierte en estructuras de datos internas del kernel:

\begin{lstlisting}[caption={Configuración estática de RTEMS en init.c}]
CONFIGURE_APPLICATION_NEEDS_CLOCK_DRIVER    // rtems_task_wake_after()
CONFIGURE_APPLICATION_NEEDS_CONSOLE_DRIVER // stdin/stdout (printf, fgets)
CONFIGURE_MICROSECONDS_PER_TICK 10000       // Tick = 10 ms (100 Hz)
CONFIGURE_UNLIMITED_OBJECTS                 // Sin limite de tareas/semaforos
CONFIGURE_UNIFIED_WORK_AREAS               // Un unico pool de memoria
CONFIGURE_INIT_TASK_STACK_SIZE (64 KB)    // Stack grande para Init
\end{lstlisting}

\subsubsection{Aplicación de testing --- \texttt{main.c}}

La aplicación tiene tres secciones claramente separadas:

\textbf{Callback de recepción \texttt{on\_rx\_data()}.} Se registra en el driver y se invoca desde la
tarea \textit{worker} cuando llegan datos. Implementa un ensamblador de líneas por canal: por cada
byte recibido, acumula hasta encontrar un carácter de fin de línea (\verb|\n| o \verb|\r|),
momento en el que imprime la línea completa con el formato
\texttt{[RX UART 02]: mensaje}. Hay un \texttt{AppLineBuffer} (1024 bytes) por cada uno de los
14 canales, declarados estáticamente en el array \texttt{rx\_lines[MAX\_TRANSCEIVERS]}, para
evitar mezclar caracteres de UARTs distintas en la misma línea de consola.

\textbf{Tarea de consola TX \texttt{Tx\_Console\_Task()}.} Es una tarea RTEMS independiente
(prioridad 100, baja) que bloquea en \texttt{fgets(stdin)} esperando comandos del usuario por
la consola USB. El protocolo de comandos es:

\begin{lstlisting}[caption={Protocolo de comandos de la aplicación de testing}]
<ID> <MENSAJE>     -> Envia MENSAJE por la UART numero ID
ALL <MENSAJE>      -> Envia MENSAJE por TODAS las UARTs
<ID> SLO ON        -> Reconfigura UART ID con Slew Rate limitado
<ID> SLO OFF       -> Reconfigura UART ID con Slew Rate normal
ALL SLO ON/OFF     -> Aplica configuracion SLO a todas las UARTs
\end{lstlisting}

Los comandos SLO re-ejecutan \texttt{Transceiver\_Init()} con una configuración distinta, lo que
permite reconfigurar el hardware en caliente sin reiniciar el sistema.

\textbf{Función \texttt{Init()} --- punto de entrada RTEMS.} RTEMS la llama al arrancar en lugar
de \texttt{main()}. La secuencia de inicialización es:

\begin{enumerate}
    \item \texttt{mmu\_map\_pl\_axi\_early()}: mapea en la MMU la región del PL para acceso \textit{bare-metal}.
    \item \texttt{Transceiver\_Global\_INIT()}: descubre el hardware e instala la ISR maestra.
    \item Bucle de inicialización: \texttt{Transceiver\_Init()}, \texttt{Transceiver\_SetRxCallback()} y
    \texttt{Transceiver\_SendString("UART N Lista.\textbackslash r\textbackslash n")} para cada UART.
    \item Creación y lanzamiento de \texttt{Tx\_Console\_Task}.
    \item \texttt{rtems\_task\_delete(RTEMS\_SELF)}: la tarea Init se elimina a sí misma; el
    kernel queda con las \textit{worker tasks} + la tarea de consola.
\end{enumerate}

\subsection{Aplicación de testing de las placas CDHS y AOCS}

\subsubsection{Ampliación en PL para soportar interfaces adicionales CDHS}

Para poder probar la funcionalidad de la placa CDHS, fue necesario añadir soporte firmware
que controlara las líneas adicionales a RS422 y RS485. La configuración del entorno fue la
siguiente: se preparó el hardware en la PL para probar las líneas de RS422/RS485, SPI, CAN
y PWM creando un proyecto en Vivado con el script TCL \texttt{regenerate\_all} con 3
transceptores serie, y añadiendo el bloque de control PWM. Además se configuró el PS para
conectar SPI y CAN con el exterior activando las líneas de SPI 0, CAN 0 y CAN 1 como
interfaces externas hacia el conector FMC.

Para controlar las líneas de PWM se instanció el bloque \texttt{PWMx4\_auto\_test}, que genera
desde la PL cuatro señales PWM a 10~kHz, 5~kHz, 1~kHz y 100~Hz con \textit{duty cycle} del
50\%, sin necesidad de ningún driver en el PS. Se utilizó el mismo BOOT.bin para todas las
pruebas de CDHS.

Para probar el ADC ADS7950 de la placa CDHS, se desarrolló una pequeña aplicación de
lectura SPI sobre RTEMS. Dado que el BSP de RTEMS 7 para ZCU102 no incluía en ese
momento un driver SPI de alto nivel, se accedió directamente al controlador SPI Cadence
integrado en el PS mediante mapeado de registros en memoria (MMIO), siguiendo el mapa de
registros descrito en el Manual de Referencia Técnico del Zynq UltraScale+
\cite{xilinx2023ug1085}. El protocolo de comunicación con el ADC ---formato de trama de
16 bits, selección de canal en \textit{Manual Mode} y gestión de la latencia de conversión--- se
implementó conforme al \textit{datasheet} del ADS7950 \cite{ti2018ads7950}.

Para probar el CAN, se utilizó una aplicación desarrollada por Diego Ramos, compañero en
el laboratorio B105 y en el proyecto LINCE. En su aplicación, se configuran las líneas de CAN
y se establecen tests que comprueban que la conexión por CAN funciona correctamente.

\subsubsection{Ampliación para soportar interfaces adicionales AOCS}

Para probar la placa AOCS, se generó un proyecto de Vivado con \texttt{regenerate\_all} con
5 transceptores, añadiendo un bloque VHDL para controlar los puentes en H por PWM
(\texttt{Motor\_H\_bridge\_test.vhd}). Este bloque genera 3 PWMs y los enruta por la línea 1 o
línea 2 de cada motor en función de la dirección deseada. Para estas pruebas, se fijó un
tiempo para cada dirección de manera que durante unos segundos el PWM fuera en un
sentido (línea 1 activa, línea 2 a 0) y durante los siguientes segundos en el sentido contrario.

La comunicación por SpaceWire no se llegó a probar por falta de tiempo al adquirir los
arneses Micro-D9.

\subsection{Validación de hardware: arneses y equipamiento}

\subsubsection{Arnés para verificar comunicaciones RS422}

Arnés a 4 hilos, cruzando las líneas TX-P/N de un driver con las RX-P/N del opuesto.

\subsubsection{Arnés para verificar comunicaciones RS485}

Arnés a 2 hilos, en el que se conectan A+ y B$-$ de un driver con los del otro.

\subsubsection{Arnés para verificar comunicaciones CAN}

Se conectan las líneas CAN\_H y CAN\_L del canal nominal con las del canal redundante.

\subsubsection{Equipo de laboratorio utilizado}

Se utilizó un osciloscopio y una fuente de alimentación de hasta 32~V.
